\newpage
\begin{center}
  \textbf{\large 3. РАЗРАБОТКА ОТЕЧЕСТВЕННОГО АНАЛОГА ПРИЛОЖЕНИЯ НА ОСНОВЕ СОБРАННЫХ ДАННЫХ}
\end{center}
\refstepcounter{chapter}
\addcontentsline{toc}{chapter}{3. РАЗРАБОТКА ОТЕЧЕСТВЕННОГО АНАЛОГА ПРИЛОЖЕНИЯ НА ОСНОВЕ СОБРАННЫХ ДАННЫХ}
% 4 строки

В данной главе для создания альтернативы системы с закрытым исходным кодом используется методика анализа защищённости мобильных приложений при помощи перехвата сетевых пакетов прокси-сервером. Как уже было отмечено в предыдущей главе, данный метод позволяет успешно судить о внутренней архитектуре и структуре неизвестной системы. Теперь же требуется оценить данный метод с точки зрения разработки на основе его результатов отечественного аналога приложения или какой-либо его части.

В первую очередь проверим возможность создания незаметного прокси-сервера, постоянно анализирующего входящий и исходящий трафик в промышленных масштабах. Такое применение может потребоваться для введения цензуры, эффективного ведения информационной войны или слежки за пользователями и компаниями, для обнаружения угроз безопасности недоступных только при рассмотрении приложения с точки зрения их исходного кода.

Во-вторых, важна возможность создания альтернативного клиента, в условиях блокировок и ограничения распространения приложений на территории Российской Федерации.

В-третьих, в случае обнаружения уязвимостей в мобильном приложении, интересно рассмотреть способы временного исправления ошибок и создания альтернативных систем в серверной части мобильного приложения. И в итоге оценить количество получаемой от метода информации для копирования всего приложения целиком с открытой реализацией.

\section{Создание постоянного прокси-сервера для автоматизированного наблюдения}
% 10 строк

На примере игры Minecraft рассмотрим создание проксирующего сервера для постоянного сбора релевантной информации об игроках.

В коллаборативной среде онлайн игры постоянно происходят события, порождённые разными игроками. С точки зрения безопасности персональных данных, отправка событий с личного устройства на устройства других пользователей без должных ограничений на доступ к персональной информации приводит к утечкам данных. Поэтому так важно собирать пересылаемые по сети данные постоянно в автоматическом режиме.

Анализ сетевых пакетов позволяет целиком воссоздать схему коммуникации клиентской и серверной стороны и поставить между ними специальный проксирующий сервер, который дополнительно проводит сбор данных на предмет соответствия нормам безопасности. И в данном примере, в качестве эксперимента, собирались данные о всех игроках, посещающих сервер одновременно и их друзьях.

В ходе сбора данных можно проверять их на наличие пропаганды, фальшивых пользователей, цензуры и других методов манипуляции. Реализованное для этого приложения решение позволило проверить оригинальность данных и после продолжительного времени незаметного наблюдения обнаружить подтверждение известной "теории шести рукопожатий" о диаметре социального графа.

Таким образом, разработанное решение помогло провести анализ защищённости приложения, оценить риски утечки данных и наблюдать за состоянием системы в будущем. Это важно, так как зачастую приложения обновляются и новые версии могут содержать неизвестные прежде уязвимости, о которых можно узнать сравнивая сетевые пакеты старой и новой версии друг с другом. Постоянный мониторинг передаваемых данных повышает доверие к программному обеспечению и не требует при этом активного вмешательства в разработку конкурентного аналога или альтернативного клиента, что позволяет обеспечивать качество с небольшим вложением ресурсов, времени разработчиков и рисков.

\section{Разработка альтернативного клиента}
% 103 строки

Пассивного наблюдения за системой не всегда достаточно для решения существующих проблем в безопасности приложения и точно недостаточно для разработки отечественного аналога программного обеспечения недоступного для распространения на территории Российской Федерации. Поэтому для эксперимента по исследованию возможности разработки альтернативного клиента на основе данных отчёта о перехваченном проксирующим сервером трафика было выбрано недоступное к установке приложение для заказа еды.

После проведения анализа приложения были выбраны ключевые функции, необходимые для работоспособности альтернативного клиента. К этим функциям относится просмотр меню, получение сведений о скидках и предложениях на основе собранной корзины, создание заказа, проведение оплаты и проверка выдачи заказа пользователю. Именно эти функции обеспечивают главное конкурентное преимущество приложения для заказа еды по сравнению с традиционным заказом с участием человека, поэтому их воспроизводимость является критерием успешного применения метода для копирования главного функционала приложения.

При более глубоком анализе системы, обеспечивающей просмотр меню, был выявлен недостаток оригинального клиента, заключающийся в отсутствии программного способа создавать новые позиции. Такой архитектурный подход обязывает обновлять приложение для каждого обновления меню, то есть каждый сезон, что создаёт ещё больше нагрузки на разработчиков оригинальной системы и плодит ошибки в системе. Отсутствие гибкости в главной функции приложения требует способа узнавать актуальный статус меню другим способом. Для этого исходное приложение пересылает на закрытые сервера компании персонально идентифицирующие данные пользователя и в том числе, его местоположение в момент совершения заказа. Лишние данные в клиентском приложении повышают размер устанавливаемых ресурсов и затрудняют автоматический анализ на уязвимости при помощи антивирусного программного обеспечения. Такое решение ставит безопасность пользователей под угрозу и требует замены уязвимого кода на безопасное решение. Конкретная реализация замены уязвимого модуля будет рассмотрена в следущей секции.

После исследования части приложения, ответственной за оценку стоимости корзины и получение данных о специальных предложениях и скидках, была обнаружена передающаяся с сервера излишняя информация о внутреннем устройстве системы, что создаёт риск атаки и вывода из строя целиком системы по оценке заказа. Такая информация не должна покидать закрытый серверный контур, так как её утечка создаёт угрозу функционированию приложения. Для обеспечения безопасности внешней системы требуется разработать дополнительный прокси-сервер, отфильтровывающий все критичные секретные данные.

Также при успешном проведении оплаты и подтверждении от платёжной системы на клиентскую часть приложения поступают секретные коды платёжной системы. Сами по себе они не представляют угрозу безопасности или компрометации пользовательских данных, но создают дополнительную поверхность для атаки и не несут при этом никакой ценности для конечного пользователя. Разработка проксирующего сервера для платёжной информации обеспечит повышение безопасности и стабильности серверной части мобильного приложения.

В совокупности эти факторы оправдывают необходимость в разработке альтернативного клиента и выделенного проксирующего сервера для улучшения защищённости приложения с закрытым исходным кодом. Эта разработка несёт практическую ценность и подтверждает утилитарность выбранного метода в применении к реальному приложению, имеющему ценность при дистрибуции.

В качестве архитектуры альтернативного клиента подходит веб клиент, потому что переносимость между различными платформами для мобильных приложений упрощает разработку и сокращает затраты ресурсов на реализацию. Проектирование и разработка конкурирующего программного обеспечения с открытым исходным кодом, решающего проблемы безопасности и риски утечки персональных данных, создаёт дополнительную ценность и позволяет мобильному приложению работать без ограничений на территории Российской Федерации.

Проведённый в рамках работы эксперимент подтвердил правильность выбора метода и продемонстрировал качественные результаты при тестировании. Разработанный альтернативный клиент открывает безопасный доступ к приложению и показывает практическую ценность предложенного метода.

\section{Замена уязвимого модуля закрытой системы}
% 44 строки

В ходе исследования протоколов обмена данными мобильного приложения для заказа еды был обнаружен потенциально вредоносный модуль, нуждающийся в замене. Проксирующий сервер, созданный совместно с альтернативным мобильным клиентом, позволяет встраивать или заменять модули собственной разработки или аналоги с открытым исходным кодом.

Найденный недостаток мобильного приложения заключается в необходимости загружать все данные, которые могут потребоваться в ходе исполнения программы, внутрь исполняемого файла, увеличивая его размер и затрудяя анализ при помощи автоматизированных инструментов по проверке мобильных приложений на защищённость.

Для предотвращения потенциальных рисков безопасности был выбран модуль для хранения большого количества конфигурационных данных с открытым исходным кодом и к нему обеспечен доступ через поддержку в проксирующем сервере и клиентском приложении. Помимо уменьшения размеров альтернативного клиента, новый подход также повышает удобство в разработке и поддерживании мобильного приложения и разделяет циклы разработки приложения и обновления конфигураций, требующихся для его функционирования.

В итоге, исследуемая методика позволила найти и заменить конкретный модуль мобильного приложения, вызывающий наибольшие риски. Метод предоставил исчерпывающее количество сведений о работе приложения, что позволило встроить новый модуль независимо в работу системы и перейти на более прозрачную и понятную для аудита архитектуру.

\section{Заключение главы}
% 75 строк
% 7 строк

В данной главе была рассмотрена методика анализа защищённости мобильных приложений с использованием перехвата сетевых пакетов прокси-сервером для разработки отечественного аналога приложения с закрытым исходным кодом. Проведённый анализ показал, что данный метод позволяет эффективно изучать внутреннюю архитектуру и структуру неизвестных систем, выявлять уязвимости и собирать данные, необходимые для создания альтернативных решений.

На примере игры Minecraft было продемонстрировано создание постоянного прокси-сервера, обеспечивающего автоматизированный мониторинг сетевого трафика. Это позволило оценить риски утечки персональных данных, выявить потенциальные угрозы безопасности, такие как пропаганда или манипуляции, и подтвердить гипотезы о структуре социального графа. Такой подход доказал свою эффективность для наблюдения за состоянием системы и выявления изменений в новых версиях приложений без необходимости активного вмешательства в их разработку.

Разработка альтернативного клиента на основе анализа сетевых данных, проведённая на примере приложения для заказа еды, показала возможность воспроизведения ключевых функций, таких как просмотр меню, оформление заказа и проведение оплаты. Выявленные недостатки оригинального приложения, включая отсутствие гибкости в обновлении меню и передачу избыточной информации, подчеркнули необходимость создания более безопасных и эффективных решений. Разработанный веб-клиент с открытым исходным кодом подтвердил практическую ценность метода, обеспечивая переносимость, снижение затрат на разработку и устранение рисков, связанных с ограничениями распространения приложений в Российской Федерации.

Замена уязвимого модуля в серверной части приложения продемонстрировала возможность интеграции открытых решений для повышения безопасности и оптимизации работы системы. Новый модуль для хранения конфигурационных данных уменьшил размер клиентского приложения и упростил его разработку и поддержку, сохраняя функциональность и повышая прозрачность архитектуры.

В совокупности, методика перехвата сетевых пакетов показала высокую эффективность для анализа и разработки отечественных аналогов приложений. Она обеспечивает достаточный объём информации для создания альтернативных клиентов, устранения уязвимостей и повышения безопасности систем, что делает её ценным инструментом для реализации конкурентоспособных решений с открытым исходным кодом, соответствующих требованиям локального рынка и нормам безопасности.
